\section*{Chapter 9 --- The transport seam and the register access}
\addcontentsline{toc}{section}{Chapter 9 --- The transport seam}

\solution{ex:9:1}{Predict}
\ans{a} \code{RegTxDataHi} has index 4. Write: \code{CmdWrite | 4} = \code{0x84}. Read: just the
index, that is, \code{0x04}.
\ans{b} \code{83 AA BB 00 00} --- the command byte \code{CmdWrite | 3}, then the value MSB first.
\ans{c} \code{0x12345678}. The first byte, \code{0xFF}, is the answer to the command byte and is
discarded; the four following ones are assembled.
\ans{d} You get \code{0xFF123456}: the value is shifted by one byte, the last data byte falls off,
and the junk byte ends up at the top. The symptom is register values that look $256$ times too
large with an inexplicable \code{0xFF} at the top.

\solution{ex:9:2}{The sign extension}
The exercise asks for a code change; the questions are answered here.
\ans{a} \code{value} becomes \code{0xFFFFFFFF80000000} truncated to 32 bits, that is,
\code{0x80000000} --- but the route there goes through undefined behaviour. \code{transfer()}
returns a \code{std::uint8_t}, which is promoted to \code{int}. \code{0x80} shifted 24 steps to
the left sets the sign bit in a 32-bit \code{int}, and shifting into the sign bit is undefined in
C++17. In practice the result is a negative number that is then converted to \code{std::uint32_t}
with ones where there should be none.
\ans{b} For every byte where bit 7 is \emph{low}, that is, all values below \code{0x80}. The code
therefore works perfectly for small identifiers and small DLC values.
\ans{c} The driver would have worked for \code{id} up to \code{0x7F} in the topmost data byte and
broken above that. Concretely: \code{0x123} works, \code{0xAABB} in the data field does not. At
bring-up it looks as though some frames arrive and others do not, depending on the \emph{content}
--- which is one of the most confusing symptoms there is, since you go looking for a fault in the
transfer when the fault is in a type conversion.

\solution{ex:9:3}{A transport to debug with}
The exercise asks for code. The output should be comparable line by line against the example in
\kapref{ch:spi}: five bytes between a \code{select} and a \code{deselect}, the command byte first,
the value MSB first.

Two things are worth building in from the start: print \emph{both} the sent and the received byte
on every line, so that a read can be followed, and count the transactions, so that a \code{send()}
can be checked against the six it should cost.

\solution{ex:9:4}{The layer boundary}
\ans{a} No. \code{ByteTransport} moves bytes and knows nothing about frames. A dependency on
\code{frame.h} would be the first step towards it knowing what it transports.
\ans{b} No. Then it would no longer be a transport but a register access, and that belongs in
\code{Spi} --- the only layer that is to know the register map.
\ans{c} No, and this is the one that sounds the most reasonable. The argument for is that the
transaction is always five bytes anyway, so why not let the transport handle it? The argument
against is that the \code{SS} framing and the five-byte format belong to the \emph{protocol}, not
to the transfer. Put them in the transport and every transport implementation --- including the
scripted test double --- has to know the protocol, and then there is no seam left: you can no
longer swap out how bytes are moved without swapping out what they mean at the same time.
\ans{d} No. Testing a private method directly is almost always a sign that it ought to be public,
or that the test is testing the wrong thing. Here it is tested through \code{send()} and
\code{receive()}, and by inspecting which bytes came out --- which is precisely what has to be
true.
